%%%%%%%%%%%%%%%%%%%%%%%%%%%%%%%%%%%%%%%%%%%%%%%%%%%%%%%%%%%%%%%%%%%%%%%%%%%%%%%%%%%%%%%%%%%%%%%%%%%%%%
%
%   Filename    : chapter_6.tex 
%
%   Description : This file will contain details of the conclusion and the recommendations of the study.
%                 
%%%%%%%%%%%%%%%%%%%%%%%%%%%%%%%%%%%%%%%%%%%%%%%%%%%%%%%%%%%%%%%%%%%%%%%%%%%%%%%%%%%%%%%%%%%%%%%%%%%%%%

\chapter{Conclusion and Recommendation}

The WeCyclePH project set out to explore the intersection of crowdsourcing, mobile technology, and urban mobility, with the objective of empowering cyclists to actively contribute to the assessment and improvement of cycling infrastructure in the Philippines. Through the development of a mobile application, the implementation of a structured data pipeline, and the deployment of a fine-tuned instance segmentation model, the project successfully demonstrated the viability of a cyclist-centered data collection platform. By capturing both objective geospatial data and subjective user annotations, WeCyclePH generated a multifaceted dataset capable of informing safer, more accessible urban transport planning.

The success of the platform is reflected in its ability to collect rich, context-aware data directly from the lived experiences of cyclists. The application's core functionality, ride recording, predefined annotation tagging, and video-enabled reporting, lowered the barrier to participation while ensuring meaningful engagement from users. Importantly, the feedback gathered during public beta testing revealed strong intrinsic motivations among cyclists, such as a sense of civic responsibility and community contribution. These motivations were further strengthened by the app's usability-focused design, which emphasized simplicity, familiarity, and intentional interaction over technical complexity.

Nonetheless, the project was not without limitations. Usage patterns revealed a clear bias toward urban and digitally connected cyclists, limiting the representativeness of the dataset across more diverse geographic and demographic groups. From a technical perspective, the instance segmentation model demonstrated promising results for common objects such as cars and bike lanes, but struggled with less frequent or ambiguous classes like potholes, debris, and localized obstructions. In addition, practical challenges, such as mobile upload failures, safety concerns during manual annotation, and limited discoverability of key features, highlighted areas requiring refinement.
Despite these constraints, the WeCyclePH platform offers a compelling proof of concept for civic data collection in support of cyclist safety and infrastructure development. The combination of machine intelligence and human insight enabled the creation of a novel dataset grounded in real-world usage, and the project laid the groundwork for scalable, participatory approaches to urban planning. With sustained iteration, targeted outreach, and further technical improvements, WeCyclePH has the potential to evolve into a foundational tool for cities seeking to prioritize active transportation and human-centered design.

\section{Recommendations for Future Work}

Building on the initial success and critical insights gained from the WeCyclePH deployment, several key directions are recommended to enhance the platform's impact, inclusivity, and technical robustness in future iterations.

First, addressing sampling bias must be a priority to ensure that the dataset reflects the diversity of the broader cycling population. The current user base was largely composed of urban, tech-savvy cyclists, often motivated by structured incentives. Future deployments should actively target underrepresented groups, including casual riders, rural cyclists, and individuals with limited digital literacy, through outreach campaigns, simplified onboarding processes, and support for multiple local languages. Doing so would not only improve the representativeness of the collected data but also ensure that the benefits of cyclist-driven infrastructure planning are more equitably distributed.

Second, significant improvements in technical performance and platform resilience are necessary. One immediate priority is the implementation of offline-first functionality, enabling users to record and annotate rides without requiring a stable internet connection. This is especially critical for cyclists operating in low-connectivity areas or for those whose devices struggle with real-time uploads. Additionally, the current reliance on cloud-based inference introduces latency and computational bottlenecks. Exploring on-device inference models or hybrid edge-cloud architectures could enable real-time annotation support, reduce data usage, and enhance user experience.

Enhancing the annotation process itself presents another key opportunity. Participants consistently highlighted the challenge of manually annotating hazards mid-ride due to safety concerns and cognitive load. Future work should investigate hands-free annotation mechanisms, such as voice-activated input, sensor-based triggers, or GPS-based auto-prompts. Post-ride annotation surveys, simplified emotion tagging, and bulk annotation review could also reduce friction while enriching the dataset with more nuanced insights.

To further sustain user engagement, integrating gamification elements such as achievements, leaderboards, and milestone tracking is recommended. These features can help foster long-term usage behaviors by making contributions feel rewarding and socially visible. In parallel, external incentives, such as discount partnerships, event-based giveaways, or even civic rewards through collaboration with local government units, can provide practical encouragement for ongoing participation.

Finally, to scale WeCyclePH beyond a prototype, its long-term success will depend on sustained partnerships and system-level integration. Future work should explore opening parts of the dataset or platform APIs to advocacy groups, planners, and researchers, enabling cross-sectoral applications of the data. By positioning WeCyclePH not only as a utility for cyclists but also as a participatory data commons, the platform can become a catalyst for inclusive urban development.

To summarize all, while WeCyclePH has validated the feasibility of cyclist-centered data collection, realizing its full potential will require thoughtful improvements in inclusivity, technical infrastructure, user experience design, and civic integration. These recommendations aim to guide future work toward building a more scalable, equitable, and impactful platform for active mobility.

% While the WeCyclePH project achieved promising results in enabling cyclist-driven data collection and analysis, several limitations were encountered that inform the path for future development and research.

% One of the most prominent limitations was the sampling bias in data contributions. The majority of active users were tech-savvy, urban-based, and often incentivized, which limited the representativeness of the dataset. As a result, cycling experiences from casual riders, rural areas, and less-connected communities were underrepresented, constraining the platform's inclusivity and reach.

% ** Insert Data Quality

% From a technical standpoint, the system's instance segmentation model, while effective for common objects, struggled to detect smaller or less frequent classes such as potholes, debris, or localized barriers. Class imbalance in the training dataset and the limitations of cloud-based inference led to low recall and missed detections. Furthermore, reliance on stable internet connectivity posed challenges for real-time uploads, and some mobile-specific bugs (e.g., ride interruption when the screen is turned off) negatively affected usability.

% Given these limitations, future work should prioritize expanding the platform's accessibility and inclusiveness. This includes targeted outreach to rural communities, casual cyclists, and less digitally literate users, supported by simplified onboarding processes and multi-language tutorials. Incorporating offline functionality will also allow data collection in areas with poor connectivity, with automatic syncing when back online.
% Enhancing the annotation experience is another key area. Future iterations should explore hands-free annotation mechanisms such as voice commands or automatic prompts based on GPS patterns or ride context. Additionally, the introduction of post-ride survey questions, emotion tagging, or batch editing tools can encourage richer and safer data submission.

% To improve detection performance, further training data diversification and model enhancements are needed. This includes increasing annotated samples for underrepresented object classes and experimenting with attention-based architectures or on-device lightweight models for real-time feedback and annotation suggestions.
% Lastly, integrating gamification elements, such as achievements, ride milestones, and community leaderboards, along with real-world incentives (e.g., discounts, rewards, partnerships with LGUs or brands) can boost engagement and long-term retention. By fostering a more interactive and rewarding environment, WeCyclePH can evolve into a vibrant digital space for cycling advocacy, civic participation, and data-driven mobility planning.

% In summary, while the current implementation demonstrates the viability and value of cyclist-centered crowdsourcing, addressing these limitations through user-centered design, model refinement, and systemic engagement strategies will be essential for WeCyclePH to scale its impact and become a cornerstone of inclusive urban mobility infrastructure.
